\chapter{Spin: the spin-resolved Hohenberg--Kohn theorem}
\label{chap:spin}

Electrons with spin on a lattice of $M$ sites, with $N_\uparrow$ spin-up and $N_\downarrow$ spin-down
electrons. The hopping conserves spin, and the potential may act differently on the two spins:
$v=(v_\uparrow,v_\downarrow)$, as in collinear spin-density-functional theory (a magnetic field along a
fixed axis). The density is the pair of spin densities $(\rho_\uparrow,\rho_\downarrow)$.

Shifting $v_\uparrow$ and $v_\downarrow$ by two \emph{separate} constants changes no ground state, because
both particle numbers are fixed. This is the known non-uniqueness of spin-DFT (a uniform magnetic field
along the quantisation axis; Capelle--Vignale, Eschrig--Pickett). The theorem of this chapter says that,
for almost every potential, this is the only non-uniqueness.

\section{Uniqueness up to trivial shifts}

\begin{definition}[Rigidity and non-uniqueness modulo trivial shifts]
  \label{def:RigidMod}
  \lean{LatticeDFT.RigidMod, LatticeDFT.NonUVMod}
  \leanok
  \uses{def:NonUV}
  Let $T$ be a set of potentials (the trivial shifts). The occupations are \emph{rigid modulo $T$} if
  $\sum_iw_in_i(c)$ constant in $c$ forces $w\in T$. A potential $v$ is \emph{non-uniquely
  $v$-representable modulo $T$} if some ground state of $H(v)$ is a ground state of $H(v')$ with
  $v-v'\notin T$. For $T$ the constant potentials these are Definition~\ref{def:NonUV}.
\end{definition}

\begin{theorem}[Generic HK modulo trivial shifts]
  \label{thm:generic_hk_ae_mod}
  \lean{LatticeDFT.nonUVMod_imp_krylov_zero, LatticeDFT.generic_hk_mod, LatticeDFT.generic_hk_ae_mod}
  \leanok
  \uses{def:RigidMod, thm:nonUV_imp_krylov_zero, thm:generic_hk_ae}
  If the occupations are rigid modulo $T$ and (H$^*$) holds, the potentials that are non-uniquely
  $v$-representable modulo $T$ lie in a closed, nowhere dense, Lebesgue-null set.
\end{theorem}
\begin{proof}
  \leanok
  The proofs of Theorems~\ref{thm:nonUV_imp_krylov_zero} and~\ref{thm:generic_hk_ae} use rigidity only to
  conclude that a shared ground state with full support forces $v-v'$ to be trivial.
\end{proof}

\section{Two spin species}

\begin{definition}[Spinful configurations and Hamiltonian]
  \label{def:H0S}
  \lean{LatticeDFT.occS, LatticeDFT.kronSum, LatticeDFT.H0S}
  \leanok
  \uses{def:Conf, def:H0Fw}
  A configuration is a pair $(x_\uparrow,x_\downarrow)$ of an $N_\uparrow$- and an $N_\downarrow$-particle
  configuration. The potential has components $v_{i\uparrow},v_{i\downarrow}$, coupled to the occupations
  $n_{i\sigma}(x_\uparrow,x_\downarrow)=n_i(x_\sigma)$. The kinetic part is the Kronecker sum
  $\Gamma(-t)\otimes1+1\otimes\Gamma(-t)$, and the interaction is any diagonal function of
  $(x_\uparrow,x_\downarrow)$; the Hubbard interaction $U\sum_in_{i\uparrow}n_{i\downarrow}$ is one example.
\end{definition}

\begin{theorem}[Rigidity for two species]
  \label{thm:rigid_occS}
  \lean{LatticeDFT.rigid_occS}
  \leanok
  \uses{def:H0S, def:RigidMod, thm:rigid_occF}
  For $0<N_\sigma<M$, the spin occupations are rigid modulo the shifts $(k_\uparrow,k_\downarrow)$, constant
  on each spin.
\end{theorem}
\begin{proof}
  \leanok
  Vary $x_\uparrow$ with $x_\downarrow$ fixed and use Theorem~\ref{thm:rigid_occF}; then the same for
  $\downarrow$.
\end{proof}

\section{The leading coefficient factorises}

\begin{lemma}[Hop distance for two species]
  \label{lem:lip_kronSum}
  \lean{LatticeDFT.lip_kronSum}
  \leanok
  \uses{def:Hl, def:H0S}
  If $d_\uparrow,d_\downarrow$ grow by at most one along the off-diagonal entries of $A_\uparrow,A_\downarrow$,
  then $d_\uparrow+d_\downarrow$ grows by at most one along those of the Kronecker sum.
\end{lemma}
\begin{proof}
  \leanok
  An off-diagonal entry of the Kronecker sum moves only one of the two species.
\end{proof}

\begin{theorem}[Top coefficients of a Kronecker sum]
  \label{thm:T_kronSum}
  \lean{LatticeDFT.layerMat, LatticeDFT.T_eq_layer_pow, LatticeDFT.kronSum_pow_apply, LatticeDFT.T_kronSum}
  \leanok
  \uses{def:T, lem:lip_kronSum}
  With $w(y_\uparrow,y_\downarrow)=w_\uparrow(y_\uparrow)+w_\downarrow(y_\downarrow)$,
  \[
    T_k(y_\uparrow,y_\downarrow)=\sum_{a=0}^k\binom ka\,T^\uparrow_a(y_\uparrow)\,T^\downarrow_{k-a}(y_\downarrow).
  \]
\end{theorem}
\begin{proof}
  \leanok
  The recursion of $T$ says that $T_k(y)=(B^k)_{yc}$ for the \emph{layer matrix} $B$: $w$ on the
  diagonal, plus the entries of $A$ that increase $d$ by one. The layer matrix of the Kronecker sum is the
  Kronecker sum $B_\uparrow\otimes1+1\otimes B_\downarrow$ of the layer matrices. The two terms commute,
  so the binomial theorem gives
  $(B^k)_{yc}=\sum_a\binom ka(B_\uparrow^a)_{y_\uparrow c_\uparrow}(B_\downarrow^{k-a})_{y_\downarrow c_\downarrow}$.
\end{proof}

\begin{theorem}[The leading coefficient factorises]
  \label{thm:alpha_kronSum}
  \lean{LatticeDFT.alpha_kronSum}
  \leanok
  \uses{thm:T_kronSum, thm:power_sum, thm:ps_unique}
  If $w_\uparrow$, $w_\downarrow$ and $w$ are injective, then
  $\alpha_{(y_\uparrow,y_\downarrow)}(y_\uparrow,y_\downarrow)=\alpha^\uparrow_{y_\uparrow}(y_\uparrow)\,
  \alpha^\downarrow_{y_\downarrow}(y_\downarrow)$.
\end{theorem}
\begin{proof}
  \leanok
  By Theorem~\ref{thm:T_kronSum} and the binomial theorem,
  $T_k(y)=\sum_{z_\uparrow,z_\downarrow}\alpha^\uparrow_{y_\uparrow}(z_\uparrow)\alpha^\downarrow_{y_\downarrow}(z_\downarrow)
  \,(w_\uparrow(z_\uparrow)+w_\downarrow(z_\downarrow))^k$ for all $k$. Compare with
  $T_k(y)=\sum_z\alpha_y(z)w(z)^k$ at the node $w(y)$ (Theorem~\ref{thm:ps_unique}); by injectivity only
  $z=y$ contributes on both sides.
\end{proof}

Physically: for a strong potential the two species do not talk to each other through the hopping, so the
amplitude for moving both of them is the product of the two separate amplitudes.

\section{The theorem}

\begin{lemma}[Generic single-species potentials]
  \label{lem:exists_good_u}
  \lean{LatticeDFT.exists_good_u}
  \leanok
  \uses{thm:exists_common_ne, thm:exists_DExp_ne, lem:forms, thm:alpha_eq_D, thm:alphaRep_eq_alphaExp}
  For spinless fermions with hopping of a common sign, and any finite family of nonzero polynomials $R_k$,
  there is $u$ with injective configuration energies, all $\alpha_y(y)\ne0$, and all $R_k(u)\ne0$.
\end{lemma}
\begin{proof}
  \leanok
  The proof of Theorem~\ref{thm:hstar_lattice_fermions}, with the $R_k$ added to the polynomial conditions.
\end{proof}

\begin{theorem}[(H$^*$) for two spin species]
  \label{thm:hstar_spin}
  \lean{LatticeDFT.hstar_spin}
  \leanok
  \uses{lem:exists_good_u, thm:alpha_kronSum, lem:lip_kronSum, thm:hstar_of_alpha, def:H0S}
  On a connected graph with hopping of a common sign and any diagonal interaction, every Krylov polynomial of
  the two-species system is nonzero.
\end{theorem}
\begin{proof}
  \leanok
  Choose $u_\uparrow$ by Lemma~\ref{lem:exists_good_u}. Choose $u_\downarrow$ by the same lemma, adding the
  nonzero linear conditions
  $w_\downarrow(y_\downarrow)-w_\downarrow(z_\downarrow)\ne w_\uparrow(z_\uparrow)-w_\uparrow(y_\uparrow)$ for
  $y_\downarrow\ne z_\downarrow$, so that the combined energies are injective. By
  Theorem~\ref{thm:alpha_kronSum} every $\alpha_y(y)\ne0$; apply Theorem~\ref{thm:hstar_of_alpha}.
\end{proof}

\begin{theorem}[Spin-resolved generic Hohenberg--Kohn theorem]
  \label{thm:generic_hk_spin}
  \lean{LatticeDFT.generic_hk_spin}
  \leanok
  \uses{thm:generic_hk_ae_mod, thm:rigid_occS, thm:hstar_spin}
  Take $N_\uparrow$ spin-up and $N_\downarrow$ spin-down electrons ($0<N_\sigma<M$) on a connected graph,
  with symmetric hopping amplitudes of a common sign and any diagonal interaction (e.g.\ Hubbard). Then the
  spin-dependent potentials $(v_\uparrow,v_\downarrow)$ at which some ground state is shared with a potential
  differing by more than separate constants on the two spins lie in a closed, nowhere dense,
  Lebesgue-null set.
\end{theorem}
\begin{proof}
  \leanok
  Theorem~\ref{thm:generic_hk_ae_mod} with Theorems~\ref{thm:rigid_occS} and~\ref{thm:hstar_spin}.
\end{proof}

\section{Spin-DFT language: scalar potential and magnetic field}

Write $v_\uparrow=v-B$, $v_\downarrow=v+B$: a scalar potential $v$ and a collinear magnetic field $B$
(Zeeman coupling $-B(n_\uparrow-n_\downarrow)$, $\mu_B=1$).

\begin{theorem}[A uniform field is not determined by the ground state]
  \label{thm:groundState_field_shift}
  \lean{LatticeDFT.isGroundState_add_const, LatticeDFT.groundState_field_shift}
  \leanok
  \uses{def:H0S, def:IsGroundState}
  A potential change that is constant on every configuration does not change the ground states. In
  particular, adding a uniform constant to $v$ and a uniform constant to $B$ leaves every ground state a
  ground state.
\end{theorem}
\begin{proof}
  \leanok
  The Hamiltonian changes by a multiple of the identity, $a(N_\uparrow+N_\downarrow)+b(N_\downarrow-N_\uparrow)$.
\end{proof}

\begin{theorem}[Generic Hohenberg--Kohn theorem in spin-DFT language]
  \label{thm:generic_hk_spin_field}
  \lean{LatticeDFT.zeeman, LatticeDFT.zeemanEquiv, LatticeDFT.NonUVField, LatticeDFT.generic_hk_spin_field}
  \leanok
  \uses{thm:generic_hk_spin, thm:groundState_field_shift}
  Under the hypotheses of Theorem~\ref{thm:generic_hk_spin}, for almost every $(v,B)$ a ground state
  shared with $(v',B')$ forces both $v-v'$ and $B-B'$ to be constant. The exceptional $(v,B)$ lie in a
  closed, nowhere dense, Lebesgue-null set.
\end{theorem}
\begin{proof}
  \leanok
  $(v,B)\mapsto(v-B,v+B)$ is an invertible linear map, so preimages of closed, nowhere dense and null sets
  are closed, nowhere dense and null. Separate constants on the two spins correspond exactly to constants
  in $v$ and in $B$.
\end{proof}

So in spin-DFT the magnetic field is never unique (a uniform field can always be added), and generically
this is the only freedom.

\section{The total density: generic HK with a symmetry}

With a spin-\emph{independent} potential (the ordinary total density), the potentials $v_\uparrow=v_\downarrow$
form a null set in the space of all $(v_\uparrow,v_\downarrow)$, so the theorems above say nothing. Total spin is
conserved for every potential: eigenstates of spin $S$ vanish on some configurations for every $v$ (a
triplet never puts both electrons on one site), so the Krylov condition (H$^*$) fails. It has to be
applied sector by sector.

\begin{definition}[Sectors]
  \label{def:RigidOn}
  \lean{LatticeDFT.RigidOn, LatticeDFT.NoGS}
  \leanok
  \uses{def:NonUV}
  The symmetry is a family of orthogonal projections $P_s$ (symmetric, idempotent, $\sum_sP_s=1$), each
  commuting with $H_0$ and with every $\mathrm{diag}(n_i)$, hence with $H(v)$ for every $v$. A sector is
  \emph{rigid} if $\sum_iw_in_i(c)$ constant on its support $\{c:(P_s)_{cc}\ne0\}$ forces $w$ constant. It
  \emph{never carries a ground state} if $P_s\psi=0$ for every ground state $\psi$ of every $H(v)$.
\end{definition}

\begin{definition}[Sector Krylov polynomial]
  \label{def:symKrylovPoly}
  \lean{LatticeDFT.symKry, LatticeDFT.symKrylovPoly, LatticeDFT.eval_symKrylovPoly}
  \leanok
  \uses{def:RigidOn, def:krylovPoly}
  Let the columns of $B_s$ span sector $s$, $P_s=B_sG_sB_s^T$, with $d_s$ columns. For a configuration $c$, the
  sector Krylov polynomial is $\det(B_s^TK_c)$, where the columns of $K_c$ are $H(v)^kP_se_c$, $k<d_s$.
\end{definition}

\begin{lemma}
  \label{lem:det_symKry_eq_zero}
  \lean{LatticeDFT.det_symKry_eq_zero, LatticeDFT.hamN_comm}
  \leanok
  \uses{def:symKrylovPoly}
  If a nonzero $\varphi$ with $P_s\varphi=\varphi$ is orthogonal to every column of $K_c$, then
  $\det(B_s^TK_c)=0$.
\end{lemma}
\begin{proof}
  \leanok
  $z=G_sB_s^T\varphi$ satisfies $B_sz=P_s\varphi=\varphi\ne0$, so $z\ne0$, and
  $z^TB_s^TK_c=\varphi^TK_c=0$.
\end{proof}

\begin{theorem}[Generic HK with a symmetry]
  \label{thm:generic_hk_sym}
  \lean{LatticeDFT.nonUV_sym_imp, LatticeDFT.generic_hk_sym}
  \leanok
  \uses{lem:det_symKry_eq_zero, def:RigidOn, thm:zeroSet_closed_nowhere, thm:volume_zeroSet}
  Assume every sector is rigid or never carries a ground state, and every sector Krylov polynomial (for
  $c$ in the support of the sector) is nonzero. Then the non-uniquely $v$-representable potentials lie in a
  closed, nowhere dense, Lebesgue-null set.
\end{theorem}
\begin{proof}
  \leanok
  A shared ground state $\psi$ has a nonzero sector component $\varphi=P_s\psi$, again an eigenvector of both
  Hamiltonians. That sector carries a ground state, so it is rigid. If $\varphi$ is nonzero on the whole
  support, rigidity forces $v-v'$ constant. Otherwise $\varphi_c=0$ for some $c$ in the support; then
  $\varphi^TH^kP_se_c=E^k\varphi_c=0$, and the sector Krylov polynomial vanishes at $v$.
\end{proof}

\begin{theorem}[Small Hubbard models, exact certificates]
  \label{thm:hubbard_certificates}
  \lean{LatticeDFT.noGS_D, LatticeDFT.hubbard_dimer_generic_hk_ae, LatticeDFT.hubbard_triangle_generic_hk_ae}
  \leanok
  \uses{thm:generic_hk_sym}
  One spin-up and one spin-down electron, $t=1$, $U=4$, spin-independent potential. On the dimer (half
  filling) and on the triangle, the total density determines the potential up to a constant for almost
  every potential. On the dimer, the triplet (density $(1,1)$ for every $v$) is not rigid, but it is never a
  ground state: its energy is $v_0+v_1$, while the trial state $(1,2,2,1)$ has energy $v_0+v_1-\tfrac45$.
\end{theorem}
\begin{proof}
  \leanok
  Singlet and triplet sectors with integer bases. Every sector Krylov determinant is evaluated exactly at
  one integer potential (generated by \texttt{scripts/gen\_hubbard\_sym\_lean.py}); rigidity from a few
  configurations.
\end{proof}

For general lattices the sector Krylov condition is open: in the strong-potential limit the spin
configurations of each occupation pattern stay degenerate, and are split only at higher order
(superexchange).
